% =========================================================================
% 七、问题三：滚动合同调整模型
% =========================================================================

\section{问题三：多时刻预报下滚动合同调整模型的建立与求解}

\subsection{信息结构与结算规则}
\label{sec:q3-frame}

问题二只有``计划—执行''两层：0:00 计划一经制定即为结算依据，实测缺口只能按 5 倍价紧急补足。问题三在两点上改变了这一结构：

\begin{enumerate}
  \item \textbf{信息增加}：除 0:00 外，6:00、12:00、18:00 还可获得未来 24 小时整点光伏预报（附件 3），因而日内可以滚动更新对剩余时段光伏出力（进而净负载）的认识；
  \item \textbf{结算增加}：允许在预报时刻\textbf{重新约定}剩余时段的购电量，但双向偏离 0:00 计划都要付费——\textbf{计划高于调整}的差额按 $0.5p_t$ 的违约电价结算，\textbf{调整高于计划}的差额按 $1.5p_t$ 结算；紧急购电仍为 $5p_t$。
\end{enumerate}

由此形成``计划—调整—执行''三层结构：在问题二框架（图~\ref{fig:q2-frame}）的``计划层''与``执行层''之间，插入一个在每个预报时刻重复触发的\textbf{调整层}。此时需定量回答两个问题：\textbf{(i)} 日内调整机会本身的价值；\textbf{(ii)} 在 6:00、12:00、18:00 三个时刻使用预报（而非仅在 0:00 使用一次）是否经济。后者是题面明确要求分析的问题。

\subsection{预报通道与净负载预测}
\label{sec:q3-forecast}

\textbf{（一）整点预报转区间电量。}附件 3 只给出逐日四个发布时刻对未来 24 小时\textbf{整点}的光伏功率预报。记发布时刻为 $\tau\in\{0,6,12,18\}$、其预报曲线为 $f_\tau(h)$，将其在相邻整点间线性插值、再在 10 分钟区间上积分，得区间电量
\begin{equation}
  \widetilde V_t(\tau)=\tfrac12\bigl[f_\tau(\tau_t)+f_\tau(\tau_{t+1})\bigr]\Delta t ,
  \label{eq:q3-pv}
\end{equation}
其中 $\tau_t=t\Delta t$ 为区间起点时刻。$\tau>0$ 时曲线端点以\textbf{当日已实测光伏功率的水平}锚定（而非初值或瞬时值），以削弱插值端点的噪声；$\tau=0$ 通道该时刻尚无实测，端点取零。

\textbf{（二）0:00 通道的预报融合。}问题二预测器所用的光伏分量来自历史经验形状，而附件 3 提供了当日的光伏预报，故可对预测作一步收缩修正：
\begin{equation}
  \widehat N_t \;=\; \underbrace{\bigl(m^{\mathrm{r}}+\gamma\Delta^{\mathrm{w}}+M_\beta\bigr)_t}_{\text{问题二预测器（式~\eqref{eq:q2-forecast}）}}\;-\;\kappa_g\Bigl[\widetilde V_t(0)-\widehat V_t^{\mathrm{hist}}\Bigr],
  \label{eq:q3-fuse}
\end{equation}
$\widehat V_t^{\mathrm{hist}}$ 为问题二预测器隐含的光伏分量，$\kappa_g\in[0,1]$ 为\textbf{分档收缩系数}，按预报时长 $h$ 分三档 $g\in\{[0,6),[6,12),[12,24)\}$。$\kappa_g$ 由\textbf{过去 $K=30$ 天}的因果滚动回归估计：
\begin{equation}
  \kappa_g=\mathrm{clip}\!\left(-\frac{\sum_s r_s x_s}{\sum_s x_s^{2}},\,0,\,1\right),\qquad
  r_s=N^{(s)}-\widehat N^{(s)}_{\text{基准}},\quad x_s=\widetilde V^{(s)}(0)-\widehat V^{(s)}_{\mathrm{hist}} ,
  \label{eq:q3-kappa}
\end{equation}
即用``实测净负载残差''对``两条预报通道的差值''回归；取负号是因为光伏预报上调（$x>0$）意味着净负载下调，截断在 $[0,1]$ 以避免过度修正。样本不足 7 天时取 $\kappa_g=0$，退回问题二预测器。$\kappa_g$ 每日仅用 $d-K,\dots,d-1$ 的数据重估，保持与问题二一致的因果口径。

\subsection{调整层：带双向违约价的剩余时段线性规划}
\label{sec:q3-adjust}

在 $\tau\in\{6,12,18\}$，先按\textbf{当前生效合同}执行完上一段区间，得到该时刻的实测储电量 $E_\tau$；再对剩余区间 $t\ge t_\tau$ 重新求解合同 $b'_t$：
\begin{equation}
\begin{aligned}
  \min\;\; & \sum_{t\ge t_\tau} p_t b'_t \;+\; \tfrac12\!\!\sum_{t\ge t_\tau}\! p_t\,(u_t+v_t) \\
  \text{s.t.}\;\;
  & u_t\ \ge\ b_t^{\mathrm{plan}}-b'_t,\qquad v_t\ \ge\ b'_t-b_t^{\mathrm{plan}},\qquad u_t,\,v_t\ge 0, \\
  & b'_t+d_t\ \ge\ \widehat N_t(\tau)+c_t,\qquad E_{t+1}=E_t+\eta_c c_t-\frac{d_t}{\eta_d},\\
  & E_{\min}\le E_t\le E_{\max},\qquad 0\le c_t,\,d_t\le \tfrac{5000}{6},\qquad b'_t\ge0,\\
  & E_{t_\tau}=E_\tau\ (\text{实测}),\qquad E_H=E_0^{(d)} ,
\end{aligned}
  \label{eq:q3-adjust}
\end{equation}
其中偏差始终对\textbf{0:00 计划} $b^{\mathrm{plan}}$ 计量（而非对上一次合同），故连续两次调整不产生重复计费。末式为计划层的日循环边界，同 \S\ref{sec:q2-lp}。

\textbf{目标函数与题面结算规则的等价性。}式~\eqref{eq:q3-adjust} 只需一次线性规划即可同时表示下调的违约费和上调的加价，无需引入分段逻辑：
\begin{itemize}
  \item 当 $b'_t\le b_t^{\mathrm{plan}}$ 时，题面结算为 $p_tb'_t+\tfrac12p_t\bigl(b_t^{\mathrm{plan}}-b'_t\bigr)$，而 $\tfrac12p_t(u_t+v_t)$ 在该情形下恰取 $u_t=b_t^{\mathrm{plan}}-b'_t$、$v_t=0$，二者完全一致；
  \item 当 $b'_t\ge b_t^{\mathrm{plan}}$ 时，$p_tb'_t+\tfrac12p_t\bigl(b'_t-b_t^{\mathrm{plan}}\bigr)=\tfrac32p_tb'_t-\tfrac12p_tb_t^{\mathrm{plan}}$，正是``超出部分按 1.5 倍、未超出部分按 1 倍''的结果。
\end{itemize}
\textbf{调整层必须自足，不含紧急购电项。}供给约束要求合同与储能放电覆盖预测需求（$b'_t+d_t\ge\widehat N_t(\tau)+c_t$），超出部分弃置。若在调整层中引入紧急购电并按其 5 倍价计入目标，线性规划会把缺口``计划''给紧急购电，从而低估所需的合同量、把风险推给执行层；因此本文的调整层不把紧急购电当作可计划的资源。紧急购电只在执行层按实测缺口发生，费用按 \S\ref{sec:q2-exec} 以 $5p_t$ 单独核算。

\textbf{执行层。}仍采用 \S\ref{sec:q2-exec} 的尽限充放电规则，但以\textbf{最新生效合同}为交付基准；段与段之间实测储电量连续传递，日末电量结转为次日日初电量。评价期、起评储电量（10800 kWh）与问题二完全一致，便于横向比较。

\subsection{两维对照：调整机会与 0:00 预报通道}
\label{sec:q3-matrix}

为把``信息''与``调整机会''两项改进分开计量，构造 $2\times2$ 方案矩阵（图~\ref{fig:q3-matrix}）：是否使用附件 3 的 0:00 预报通道（式~\eqref{eq:q3-fuse}）$\times$ 是否允许日内三次调整。

\begin{table}[H]
  \centering
  \caption{0:00 预报通道 $\times$ 日内调整机会的费用矩阵（评价期 2025.2.1--12.31，334 天，万元）}
  \label{tab:q3-matrix}
  \begin{tabularx}{\textwidth}{lXcc}
    \toprule
    方案 & 设定 & 总费用 & 相对 A0 \\
    \midrule
    A0 & 不用 0:00 预报通道，不允许日内调整（$=$ 问题二主方案 E） & 1395.70 & — \\
    A0K & 用 0:00 预报通道，不允许日内调整 & 1376.20 & $-19.50$ \\
    B0 & 不用 0:00 预报通道，允许日内调整 & 1355.58 & $-40.12$ \\
    \textbf{B1（主方案）} & \textbf{用 0:00 预报通道，允许日内调整} & \textbf{1345.71} & $\mathbf{-49.99}$ \\
    \bottomrule
  \end{tabularx}
\end{table}

\begin{figure}[htbp]
  \centering
  \includegraphics[width=0.8\textwidth]{p3_fig1.pdf}
  \caption{0:00 预报通道与日内调整机会的总费用矩阵（万元）}
  \label{fig:q3-matrix}
\end{figure}

由表~\ref{tab:q3-matrix} 与图~\ref{fig:q3-matrix} 可得三点结论：

\begin{enumerate}
  \item \textbf{A0 与问题二主方案完全重合}（1395.70 万元），说明问题二、问题三两条链路在同一评价口径下自洽——问题三的``不用预报通道、不允许调整''退化为问题二模型；
  \item 单独引入日内调整降费 \textbf{40.12 万元}，单独引入 0:00 预报通道降费 \textbf{19.50 万元}；在本文的数据上，\textbf{调整机会的价值约为预报通道的两倍}；
  \item 两项单独收益之和为 59.62 万元，而联合收益为 \textbf{49.99 万元}，\textbf{重叠 9.62 万元}。这说明两类改进\textbf{部分替代}：手中已有第二次决策机会时，0:00 预报的边际价值下降。因此两个数字不能简单相加，这正是必须做 $2\times2$ 对照而非单项消融的原因。
\end{enumerate}

B1 即本文问题三主方案，全年费用 \textbf{1345.71 万元}，相对问题二主方案节省 49.99 万元（$-3.58\%$）。

\subsection{指定日期的表 1、表 2、表 3 结果}
\label{sec:q3-tables}

按题目表 1、表 2、表 3 的格式，表~\ref{tab:q3-t1}--表~\ref{tab:q3-t3} 给出四个指定日期（2025 年 3 月 20 日、6 月 21 日、9 月 23 日、12 月 21 日）的完整结果，主方案为 B1。四个日期分别对应春分、夏至、秋分、冬至，覆盖评价期内负载与光伏的季节差异。

\textbf{表 1：指定时间段购电量。}表~\ref{tab:q3-t1} 给出题目指定的六个 10 分钟时段（10:00、12:00、14:00、16:00、18:00、20:00 起点）的购电量，以及全天计划购电量、调整后购电量与全日购电费（含调整附加费与紧急购电费）。可以看到：\textbf{六个指定时段的调整后购电量与 0:00 计划一致}——调整主要发生在非指定时段，故指定时段两列数值相同；全天口径下两者出现差异（例如 12 月 21 日为计划 98349.89 kWh、调整后 96906.82 kWh）。夏季（6 月 21 日）调整后全天购电量略高于计划（36911.13 对 36864.37 kWh）：该日光伏充沛、缺口少，调整的作用是把购电在时段间微调，上调部分按 1.5 倍计价，故购电量略增而附加费仅 29.28 元。

\begin{table}[htbp]
  \centering
  \caption{四个指定日期的购电量：六个指定时段（计划值与调整后值一致）与全天口径（kWh），以及全日购电费（元）}
  \label{tab:q3-t1}
  \footnotesize
  \setlength{\tabcolsep}{3.6pt}
  \begin{tabular}{lcccccccccc}
    \toprule
    日期 & 10:00 & 12:00 & 14:00 & 16:00 & 18:00 & 20:00
         & 全天计划 & 全天调整后 & 全日购电费 \\
    \midrule
    2025-03-20 & 0.00 & 554.52 & 0.00 & 567.05 & 701.86 & 0.00
               & 68944.02 & 68455.43 & 42081.89 \\
    2025-06-21 & 0.00 & 0.00 & 0.00 & 212.16 & 0.00 & 0.00
               & 36864.37 & 36911.13 & 20752.42 \\
    2025-09-23 & 0.00 & 411.45 & 0.00 & 512.77 & 770.17 & 0.00
               & 67064.04 & 66896.01 & 45106.35 \\
    2025-12-21 & 0.00 & 1039.77 & 0.00 & 845.05 & 735.21 & 0.00
               & 98349.89 & 96906.82 & 62694.46 \\
    \bottomrule
  \end{tabular}
\end{table}

计入紧急购电后，四日的合计购电量依次为 68525.56、36911.13、67365.51、96906.82 kWh；其中 3 月 20 日与 9 月 23 日发生紧急购电（表~\ref{tab:q3-t3}），6 月 21 日与 12 月 21 日为零。费用构成的差异同样清晰：3 月 20 日合同费 41441.26 元、调整附加费 151.45 元、紧急购电费 489.19 元；9 月 23 日紧急购电费升至 2825.21 元（晚峰缺口最重）；12 月 21 日合同费 62202.88 元，为四日最高，反映冬季高负载。

\textbf{表 2：指定时间段充放电量与 0:00、24:00 储电量。}表~\ref{tab:q3-t2} 给出六个 4 小时区间的充电量与放电量，以及 0:00、24:00 的储电量。四日都呈``凌晨与午间充电、早晚峰放电、深夜回充''的形态；储电量在日间两次触及上限附近、晚峰后回落，0:00 与 24:00 的差异反映实际储电量的跨日连续衔接（如 9 月 23 日 10585.62 $\to$ 10454.35 kWh，与 \S\ref{sec:q2-days} 中问题二同日的实测轨迹一致，说明该日的合同调整未改变执行层的物理动作）。

\begin{table}[htbp]
  \centering
  \caption{四个指定日期的六段 4 小时充放电量（kWh）与 0:00、24:00 储电量（kWh）}
  \label{tab:q3-t2}
  \footnotesize
  \setlength{\tabcolsep}{3.4pt}
  \begin{tabular}{llcccccccc}
    \toprule
    日期 & 项 & 0:00--4:00 & 4:00--8:00 & 8:00--12:00 & 12:00--16:00
         & 16:00--20:00 & 20:00--24:00 & 0:00 & 24:00 \\
    \midrule
    \multirow{2}{*}{03-20} & 充 & 133.27 & 325.08 & 6061.85 & 4210.58 & 46.14 & 10717.51 & \multirow{2}{*}{10800.00} & \multirow{2}{*}{10784.00} \\
                           & 放 & 189.75 & 5470.08 & 2777.38 & 268.95 & 5803.05 & 2915.68 & & \\
    \midrule
    \multirow{2}{*}{06-21} & 充 & 169.79 & 454.45 & 9377.58 & 0.00 & 105.32 & 10048.72 & \multirow{2}{*}{10800.00} & \multirow{2}{*}{10800.00} \\
                           & 放 & 3264.87 & 4836.61 & 0.00 & 0.00 & 5738.88 & 2485.90 & & \\
    \midrule
    \multirow{2}{*}{09-23} & 充 & 330.48 & 418.77 & 5595.42 & 4850.22 & 40.07 & 10298.83 & \multirow{2}{*}{10585.62} & \multirow{2}{*}{10454.35} \\
                           & 放 & 288.54 & 6361.61 & 1896.62 & 517.88 & 6099.83 & 2396.04 & & \\
    \midrule
    \multirow{2}{*}{12-21} & 充 & 94.48 & 974.65 & 6193.26 & 6327.44 & 9.61 & 10561.20 & \multirow{2}{*}{10800.00} & \multirow{2}{*}{10800.00} \\
                           & 放 & 172.00 & 1515.29 & 6891.61 & 2428.86 & 5719.83 & 2842.53 & & \\
    \bottomrule
  \end{tabular}
\end{table}

\textbf{表 3：紧急购电结果。}表~\ref{tab:q3-t3} 按题目表 3 格式给出紧急购电的购电时间段与购电量（连续区间已合并）。四日中仅 3 月 20 日与 9 月 23 日发生紧急购电，且全部集中在晚峰时段（20:20--21:50）：此时光伏退出、净负载抬升，若合同上调幅度不足或储能已接近下限，缺口便只能按 5 倍电价补足。9 月 23 日的三个事件（共 469.50 kWh）是四日中紧急购电最重的时段。

\begin{table}[htbp]
  \centering
  \caption{四个指定日期的紧急购电结果（合并连续区间）}
  \label{tab:q3-t3}
  \small
  \begin{tabular}{llcc}
    \toprule
    日期 & 购电时间段 & 购电量/kWh & 当日合计/kWh \\
    \midrule
    2025-03-20 & 20:30--20:40 & 70.12 & 70.12 \\
    \midrule
    2025-06-21 & 无 & 0 & 0 \\
    \midrule
    \multirow{3}{*}{2025-09-23} & 20:20--20:40 & 278.48 & \multirow{3}{*}{469.50} \\
                                & 20:50--21:10 & 129.61 & \\
                                & 21:20--21:50 & 61.40 & \\
    \midrule
    2025-12-21 & 无 & 0 & 0 \\
    \bottomrule
  \end{tabular}
\end{table}

上述四个日期的结果与全年 334 天的汇总口径一致；全部逐区间策略已按题目要求写入 \texttt{result3.xlsx}（含``计划购电量''``调整购电量''``充放电量''``紧急购电量''四张表）。

\subsection{主方案的日内形态与逐月费用}
\label{sec:q3-day}

\begin{figure}[htbp]
  \centering
  \includegraphics[width=0.86\textwidth]{p3_fig3.pdf}
  \caption{2025-09-23 主方案 B1 的日内调度：(a)~实测净负载、0:00 计划合同与实际交付合同；(b)~储能电量轨迹；(c)~日内合同调整量与紧急购电量}
  \label{fig:q3-day}
\end{figure}

以 2025-09-23 为例（图~\ref{fig:q3-day}）：0:00 计划的合同 $b^{\mathrm{plan}}$（图中浅色实线）与实际交付合同 $b$（虚线）在夜间与清晨几乎重合，主要调整发生在 12:00 时刻——该次求解把午后部分区间的合同\textbf{下调}（图~\ref{fig:q3-day}c 中 12 时附近的负向柱），下调部分只需按 $0.5p$ 支付违约费；20:00 前后晚峰光伏退出、净负载抬升，合同被\textbf{上调}并伴生紧急购电（图中橙色柱）。储能轨迹仍保持``白天光伏充沛时充电、晚峰放电''的日循环，日末储电量与日初衔接，与问题二同期的运行形态一致。

\subsection{其他时刻预报的引入必要性}
\label{sec:q3-other}

这是题面明确要求回答的问题。本节考察两种``把 6:00/12:00/18:00 预报真正用起来''的做法：

\begin{itemize}
  \item \textbf{全融合（B3）}：每个调整时刻 $\tau$ 都用当期附件 3 预报重估剩余时段的净负载，即把式~\eqref{eq:q3-fuse} 从``仅 0:00 通道''推广到``当期通道''；
  \item \textbf{门控（G）}：承认``用不用''本身是一个决策，在每个决策点用纯因果的历史回放，从融合强度 $\lambda\in\{0,\tfrac14,\tfrac12,1\}$ 中择一（$\lambda=0$ 即退回 B1）。
\end{itemize}

\subsubsection{预报精度的改善}

\begin{figure}[htbp]
  \centering
  \includegraphics[width=0.95\textwidth]{p3_fig2.pdf}
  \caption{(a)~白天剩余时域中心预测 MAE（kWh/10\,min 区间）；(b)~全融合 B3 相对 B1 的费用瀑布}
  \label{fig:q3-mae}
\end{figure}

图~\ref{fig:q3-mae}a 给出白天剩余时域的中心预测平均绝对误差：全融合 B3 在 6:00 与 12:00 两个发布时刻将 MAE 由 54.34、51.04 降至 50.02、46.24，降幅分别为 4.32 与 4.80；18:00 时刻略升 0.10（28.93 $\to$ 29.03），原因是该时刻距当日光伏结束已近，可改进的时域很短。

\begin{figure}[htbp]
  \centering
  \includegraphics[width=0.95\textwidth]{p3_fig5.pdf}
  \caption{光伏区间电量 MAE：预报来源（历史预报／附件 3 预报）与端点锚定方式（历史锚点／实测锚点）的四种组合，(a)~发布后首小时；(b)~首小时之后}
  \label{fig:q3-pv}
\end{figure}

图~\ref{fig:q3-pv} 进一步在光伏区间电量本身上作四项对照。附件 3 预报叠加实测锚点的组合在两个时段上整体误差最低，且改进集中在\textbf{发布后首小时}（图~\ref{fig:q3-pv}a）——12:00 时刻的四种组合差异最大，说明当日上午实测量已经足以校正当日的光伏形状。这为式~\eqref{eq:q3-pv} 中用实测水平锚定端点的做法提供了直接依据。

\subsubsection{总费用的上升}

图~\ref{fig:q3-mae}b 给出 B3 相对 B1 的费用瀑布（表~\ref{tab:q3-b3}）：最终合同费用与紧急购电费用均略有下降，但\textbf{调整附加费大幅上升}，合计总费用上升。

\begin{table}[H]
  \centering
  \caption{全融合 B3 相对主方案 B1 的费用变化（评价期 334 天，万元）}
  \label{tab:q3-b3}
  \begin{tabular}{lcc}
    \toprule
    费用项 & 变化量 & 方向 \\
    \midrule
    最终合同费用 & $-2.39$ & 略降（更准的预测选中了更便宜的合同） \\
    调整附加费 & $+13.96$ & \textbf{显著上升（主导项）} \\
    紧急购电费 & $-1.98$ & 略降 \\
    \midrule
    合计 & $+9.59$ & \textbf{总费用上升} \\
    \bottomrule
  \end{tabular}
\end{table}

机制是明确的：精度更高的点预测使线性规划\textbf{更倾向于}偏离 0:00 计划而转向低价时段，而每一次偏离都需按 $0.5p$／$1.5p$ 付费；附加费的增量（13.96 万元）远大于合同费用与紧急费用节省之和（4.37 万元）。即\textbf{双向违约价将``精度更高的日前预测''转化为``更频繁的合同偏离''}。

\subsubsection{逐决策点自适应门控的结果}

\begin{figure}[htbp]
  \centering
  \includegraphics[width=0.96\textwidth]{p3_fig6.pdf}
  \caption{(a)~门控 $\lambda$ 在 1002 个决策点上的选择分布；(b)~门控方案 G 相对 B1 的逐月费用差（橙色表示门控方案费用更高，蓝色表示更低）}
  \label{fig:q3-gate}
\end{figure}

图~\ref{fig:q3-gate}a 给出门控 G 在全部 1002 个决策点（334 天 $\times$ 3 个调整时刻）上的选择分布：6:00 时刻 $\lambda=0$ 与 $\lambda=1$ 分别为 128 与 133 次，12:00 时刻为 124 与 134 次，18:00 时刻为 152 与 127 次，两个端点档位合计均占约七至八成，中间档位（$\lambda=0.25,0.5$）合计不足三成。这本身即表明：\textbf{在单一决策点层面难以判断融合是否有利}，预测精度与结算费用之间的方向并不总是一致，因此门控退化为在两个端点之间选择。

其结果是门控方案全年比 B1 高 \textbf{6.08 万元}（图~\ref{fig:q3-gate}b）：11 个月中仅 12 月为负（$-0.62$ 万元），2 月与 7 月增量最大（分别为 1.60 与 1.83 万元）。

\begin{figure}[htbp]
  \centering
  \includegraphics[width=0.96\textwidth]{p3_fig4.pdf}
  \caption{(a)~B1 相对 A0 的逐月节省（合计 49.99 万元，均值 4.54 万元/月）；(b)~全融合方案费用差的分组热力图（B3 $-$ B1）}
  \label{fig:q3-month}
\end{figure}

图~\ref{fig:q3-month}a 给出主方案 B1 相对基准 A0 的逐月节省：11 个月\textbf{全部为正}，合计 49.99 万元、均值 4.54 万元/月，其中 6 月最高（13.60 万元），与 \S\ref{sec:q2-year} 中 6 月紧急费用最集中相呼应——夏季晚峰缺口最需要日内的第二次决策。图~\ref{fig:q3-month}b 进一步按``电价组 $\times$ 时段''分解 B3 与 B1 的费用差：正值集中在 6--18 时的\textbf{中价组}（$+3.99$、$+3.44$），正是合同最可能被反复调整、违约附加费最易触发的时段；低价组与高价组多数取值接近零或为负，说明融合的代价并非均匀分布，而是集中在中间价位区间。

\subsubsection{小结}

在本文的结算口径与所考察的网格范围内，\textbf{不需要引入其他时刻的预报来制定调整购电策略}。附件 3 的信息\textbf{在 0:00 通道使用一次}（方案 B1，1345.71 万元）即已足够；6:00、12:00、18:00 的预报只需用于``按实测重新校正剩余时段合同''这一项工作，不必再作为新的预测输入反复重建预测。其原因不在于预报精度——附件 3 的逐时预报精度更高（图~\ref{fig:q3-mae}a、图~\ref{fig:q3-pv}）——而在于本题的双向违约价显著削弱了``预测更准''向``费用更低''的转化渠道：预测精度提高带来的合同费用与紧急费用节省（4.37 万元）不足以补偿由此增加的违约附加费（13.96 万元）。

\textbf{适用范围说明。}该结论是在``偏差对 0:00 计划计量、附加费 $0.5p$/$1.5p$、紧急价 $5p$''这套结算口径下，以及本文所考察的全融合与门控两种形式下成立的。若结算口径改变（例如允许无成本无限次调整，或偏差改为对上一次合同计量），结论需要重新核验；本文不据此断言``其他时刻预报无信息价值''。
